\section{From Local to Global: A Graph RAG Approach to Query-Focused Summarization}
\paperhead{From Local to Global: A Graph RAG Approach to Query-Focused Summarization}{Darren Edge, Ha Trinh, Newman Cheng, Joshua Bradley, Alex Chao, et al.}{arXiv preprint (Microsoft Research), 2024}{https://arxiv.org/abs/2404.16130}{arXiv API id\_list=2404.16130, title match 1.00}

\subsection{What the paper does}

Conventional RAG answers a question by retrieving the top-$k$ passages that look most like the query. That works for local questions whose answer sits in a few chunks, and fails for global questions aimed at the corpus as a whole, such as ``what are the main themes in this dataset?''. No individual chunk contains the answer, so similarity retrieval has nothing correct to return. The paper frames this as query-focused summarization rather than retrieval, and observes that classical query-focused summarization does not scale to the corpus sizes RAG systems routinely index.

GraphRAG closes that gap by moving the expensive work to index time. An LLM reads the corpus once and derives an entity knowledge graph; that graph is partitioned hierarchically into communities of closely related entities; and a summary is pregenerated for every community at every level. At query time the system never touches raw text. It runs a map-reduce over the pregenerated community summaries: each summary yields a partial answer, and the partials are combined into one global answer. The index is a precomputed, multi-resolution description of the whole corpus, so query cost is bounded by the number of community summaries rather than by corpus size.

\subsection{How the method works}

Source text is split into 600-token chunks with 100-token overlap. Chunk size is a real lever, not a detail: the authors report that GPT-4 extracted almost twice as many entity references at 600 tokens as at 2400. An extraction prompt pulls out entity instances (name, type, description), relationship instances, and covariate claims from each chunk. Because a single pass misses entities, the LLM is then asked in a self-reflection step whether anything was missed, and these ``gleanings'' loop for additional passes. Instances of the same entity or relation across chunks are then abstracted into element summaries: one description per node and per edge, with duplicate relations folded into edge weights.

The resulting weighted graph is partitioned with Leiden community detection applied recursively, giving a hierarchy from broad root communities (C0) down to leaf communities that cannot be partitioned further (C3). Each community gets a report-style summary. At leaf level, elements are added in decreasing order of combined source and target node degree until the context window fills; at higher levels, sub-community summaries are substituted for element summaries so the parent still fits. Answering is map-reduce: community summaries at a chosen level are shuffled, chunked to a fixed token size, and each chunk produces an intermediate answer plus a helpfulness score from 0 to 100. Zero-scored partials are discarded, the rest are sorted by descending helpfulness, and a reduce call writes the final answer. A fixed 8k context window was used for all query-time operations, the smallest tested and, per the paper, universally better than larger ones.

\begin{center}
\begin{tikzpicture}[
  >=Stealth, node distance=1.4cm and 0.45cm, font=\scriptsize,
  bx/.style={draw, rounded corners=1.5pt, align=center, inner sep=2.5pt, minimum height=7.5mm, text width=1.75cm},
  ar/.style={->, thick}]
\node[bx] (c1) {600-tok chunks};
\node[bx, right=of c1] (c2) {entity, relation, claim instances};
\node[bx, right=of c2] (c3) {gleaning passes};
\node[bx, right=of c3] (c4) {element summaries};
\node[bx, right=of c4] (c5) {Leiden hierarchy C0..C3};
\node[bx, right=of c5] (c6) {community summaries};
\node[bx, below=0.75cm of c3] (q1) {query};
\node[bx, right=of q1] (q2) {map: partial answer + score};
\node[bx, right=of q2] (q3) {filter 0, sort by score};
\node[bx, right=of q3] (q4) {reduce: global answer};
\draw[ar] (c1) to (c2); \draw[ar] (c2) to (c3); \draw[ar] (c3) to (c4);
\draw[ar] (c4) to (c5); \draw[ar] (c5) to (c6);
\draw[ar] (q1) to (q2); \draw[ar] (q2) to (q3); \draw[ar] (q3) to (q4);
\coordinate (mid) at ($(c6.south)+(0,-0.3)$);
\draw[ar] (c6.south) to (mid) -| (q2.north);
\end{tikzpicture}
\end{center}

\subsection{What it reports}

Two corpora were indexed: podcast transcripts (1,669 chunks, roughly 1 million tokens) and news articles (3,197 chunks, roughly 1.7 million tokens). The podcast graph held 8,564 nodes and 20,691 edges; the news graph 15,754 nodes and 19,520 edges. For each corpus, 125 global sensemaking questions were generated by having an LLM invent personas, tasks per persona, and questions per persona-task pair (5 by 5 by 5), deliberately avoiding questions answerable by a single low-level fact. An LLM judge scored head-to-head answer pairs on comprehensiveness, diversity and empowerment, with directness (how specifically an answer addresses the question) used as a validity control rather than a target. Six conditions were compared: GraphRAG at four hierarchy levels, source text summarization (TS), and a vector RAG semantic search baseline (SS).

Against the vector RAG baseline, GraphRAG won on comprehensiveness in 72 to 83\% of podcast comparisons and 72 to 80\% on news ($p<.001$), and on diversity in 75 to 82\% of podcast comparisons and 62 to 71\% on news. The cost story is on the query side: relative to summarizing the source text directly, low-level community summaries (C3) needed 26 to 33\% fewer context tokens, and root-level summaries (C0) over 97\% fewer.

\subsection{How Mytheus uses this}
\begin{takeaway}
This is the direct precedent for the Apache AGE organizational-memory graph and the graph arm of the \texttt{expand} $\to$ \texttt{retrieve} $\to$ \texttt{rerank} $\to$ \texttt{assemble} $\to$ \texttt{cite} pipeline, but adopt it offline only. Build a nightly per-tenant job that runs Leiden over the org-partitioned AGE graph (Person, Meeting, Decision, Action, Topic joined by ATTENDED, PRODUCED, ASSIGNED\_TO, RELATES\_TO, FOLLOWS\_FROM) and pregenerates one Opus community summary per cluster, stored as a fourth retrievable surface whose hits enter RRF ($k=60$) beside the Qdrant dense and Postgres full-text arms. That answers ``what has this team actually decided about billing this quarter?'', which top-10 chunk retrieval cannot. Two constraints follow. First, map-reduce global answering is unusable live: an LLM call per community-summary chunk plus a reduce call cannot fit the ContextGapAgent's 1.5\,s p95 peer round trip and would blow the p95 4\,s agent budget, so route it to the asynchronous post-meeting query surface, never the live window path. Second, community summaries are LLM abstractions with no provenance, so the \texttt{cite} stage would drop them; store each summary with the constituent \texttt{meeting\_id} and \texttt{chunk\_id} values and cite those, or the citation contract breaks. The claim to measure is index economics and precision: GraphRAG pays an extraction call per 600-token chunk plus gleaning passes plus a summary per community over an LLM-guessed entity set, exactly the naive per-chunk profile the five-lever cascade (Haiku pre-gate at roughly 35\% pass, 30-second windowing, LiteLLM tiering) removes, while NO SILENT COMMITS means Mytheus edges exist only after user confirmation. Benchmark index cost per meeting and edge precision against a GraphRAG-style extraction over the same transcripts and report the difference.
\end{takeaway}

\subsection{Limitations and what it leaves open}
\begin{itemize}[leftmargin=1.2em, itemsep=1pt, topsep=2pt]
  \item The index is built once over a corpus that is complete before indexing starts; no incremental update is evaluated, and the paper never measures the cost or latency of adding a document, which for Mytheus means a meeting ending at 11am would in principle require recomputing the Leiden hierarchy and the affected community summaries with no stated bound.
  \item Both corpora are static recorded artefacts (podcast transcripts, news articles), so the method is never tested on live streaming speech, on partial or revisable transcripts, or under any condition where the answer must land while the conversation is still open.
  \item There is exactly one corpus, one graph and one answer per question: no per-user role conditioning, no entitlement or access control on entities and community summaries, and no tenant partitioning, so the design gives no guidance on isolating one org's entities from another's during community detection.
  \item Evaluation is relative LLM-as-judge win rates on comprehensiveness, diversity and empowerment with directness as a control; there is no ground-truth answer key, no citation-resolution check, and no hallucination rate, so a more comprehensive but partly fabricated answer would score well.
  \item Query cost is reported only in context tokens, never in wall-clock seconds, and the extracted graph itself is never scored against a gold entity or relation set despite the reported sensitivity of extraction yield to chunk size, leaving both latency and index precision unknown.
\end{itemize}
